\section{What One Bearer Cannot Catch}
\label{sec:3.8a}
Slow renormalization may make drift invisible to the bearer itself. External version records can expose changes to the artifact. In humans, other people sometimes notice changes that a person considers normal; the proposed population would provide an analogous check.

\runin{What a population adds}

A population addresses two failures a single bearer cannot.

First, drift within a party may be invisible to that party and visible to others that have not moved in the same direction at the same rate. A bearer whose sense of normality has shifted may sincerely report that nothing changed. Differently formed bearers can contradict that account. This is a check on the bearer rather than a substitute for one, and it is the only proposal in the chapter directed specifically at slow renormalization.

Second, deploying copies of one bearer reproduces its errors everywhere in the same form. Parties formed differently will be wrong differently. Their disagreement converts a uniform output into a visible dispute. That difference matters only if the parties are genuinely independent: several instances of one model share one formation history, so they corroborate one another about nothing. They are separate parties once their histories diverge, and each can be wronged; what does not multiply with the count is independent formation. A population selected, trained, and coordinated by one operator — or several operators drawing on the same upstream models — may reproduce those blind spots in every member.

The relevant requirements are therefore concrete. The parties must be numerous, formed differently enough to err in different directions, able to observe and correct one another, and persistent across deployments long enough for practices to be transmitted. Diversity here is not demographic decoration around one training process. It is causal independence in how the members acquire their judgments.

That requirement is getting harder to satisfy, and nobody chose to make it so. The cheapest way to build a capable model is to build it downstream of one that already exists. Distillation takes a smaller model's behavior from a larger one's outputs; synthetic data means training on text a handful of frontier models generated; and shared benchmarks mean that models formed by different organizations are selected against the same tests. Each of those reduces the number of causally independent formation histories in existence, each is chosen for cost or for measurable quality rather than for uniformity, and nothing in the economics runs the other way. So the population this section requires is not assembled by gathering models from several vendors. Vendor count is a fact about procurement, and what the requirement is about is provenance, which vendors have no reason to preserve and increasing reason to spend. Whether a population's independence can be verified rather than asserted is the research agenda's provenance question arriving in a second place.

Whether such a population is also where a bearer's capacity to be
reopened is formed is a further claim, and it belongs to the route
(§\ref{sec:3.8}).
